\section{Cut-aware topology equivalence}
\label{app:topology-equivalence}

A cut topology is specified by
\begin{equation}
 \mathcal T=(\boldsymbol D,\boldsymbol\ell,\boldsymbol p,\mathcal K;
 C,\boldsymbol\xi,\mathcal S,\boldsymbol{\mathfrak m}),
 \label{eq:cut-topology-record}
\end{equation}
where $\boldsymbol D=(D_1,\ldots,D_N)$ is the ordered family,
$\boldsymbol\ell$ and $\boldsymbol p$ are the integration and external
momenta, $\mathcal K$ contains the external kinematic relations, $C$ is the
set of cut slots, and $\xi_c=\pm1$ is the cut energy orientation.  The
partition $\mathcal S$ distinguishes phase-space, forward-virtual, and
conjugate-virtual integration momenta.  The optional labels
$\mathfrak m_i$ record a measured or otherwise distinguished line.

Two records $\mathcal T$ and $\mathcal T'$ are equivalent when there is an
affine map
\begin{equation}
 \ell_r=\sum_sA_{rs}\ell'_s+\sum_aB_{ra}p_a,
 \qquad A_{rs},B_{ra}\in\mathbb Q,
 \qquad |\det A|=1,
 \label{eq:topology-affine-map}
\end{equation}
and a permutation $\pi\in S_N$ such that
\begin{equation}
 D_i(\ell,p)\big|_{\eqref{eq:topology-affine-map}}
 \stackrel{\mathcal K}{=}D'_{\pi(i)}(\ell',p).
 \label{eq:topology-denominator-map}
\end{equation}
No slot-dependent multiplicative factor is allowed.  The map must satisfy
\begin{equation}
 \pi(C)=C',
 \qquad
 q_i\big|_{\eqref{eq:topology-affine-map}}
 =s_iq'_{\pi(i)},
 \qquad
 \xi'_{\pi(i)}=s_i\xi_i,
 \qquad s_i=\pm1,
 \label{eq:topology-cut-map}
\end{equation}
for every cut slot.  It must also preserve the members of $\mathcal S$ and
transport each defined label $\mathfrak m_i$ to
$\mathfrak m'_{\pi(i)}$.  The determinant condition preserves the
dimensionally regulated measure.

If the source indices are $\boldsymbol\nu$, the exact change of variables
gives
\begin{equation}
 I_{\mathcal T}(\nu_1,\ldots,\nu_N)
 =I_{\mathcal T'}
 (\nu_{\pi^{-1}(1)},\ldots,\nu_{\pi^{-1}(N)}).
 \label{eq:topology-index-map}
\end{equation}
This is the powered-integral equivalence used to reuse an IBP reduction.

For example, consider
\begin{align}
 \boldsymbol D_A&=(\ell_1^2,\ell_2^2,(\ell_1-\ell_2)^2,
 (\ell_1-p)^2,(\ell_1+\ell_2-p)^2),\notag\\
 \boldsymbol D_B&=(r_1^2,r_2^2,(r_1-r_2)^2,
 (r_2-p)^2,(r_1+r_2-p)^2).
 \label{eq:topology-example-families}
\end{align}
The map $\ell_1=r_2$, $\ell_2=r_1$ has determinant $-1$ and induces
$\pi=(12)$.  If slots 1 and 2 are cuts with orientations
$(+1,-1)$ in $A$ and $(-1,+1)$ in $B$, and the remaining sector and
line labels agree, then
\begin{equation}
 I_A(\nu_1,\nu_2,\nu_3,\nu_4,\nu_5)
 =I_B(\nu_2,\nu_1,\nu_3,\nu_4,\nu_5).
 \label{eq:topology-example-identity}
\end{equation}
Keeping the orientations $(+1,-1)$ in both records would violate
Eq.~\eqref{eq:topology-cut-map}, even though the quadratic denominators still
match.
